% Independent contribution to Part A.2(b), by /root/algebra_c.
% AI-produced by GPT-6 Astra (OpenAI), ChatGPT Pro mode, 19 November 2026.
\subsection{A positive energy identity for the conformal-radius map}
\label{C:A2b-energy-section}

Fix a compact connected Riemann surface $X$ of genus $g$, distinct labelled points $p_1,\ldots,p_n$, and local holomorphic coordinates $z_i$ with $z_i(p_i)=0$. Assume $2g-2+n>0$. The coordinates specify the norms in which we measure conformal radii. For boundary lengths $\ell_i\geq0$, Mondello's infinite-grafting theorem, Theorem~5.4 (2011, p.~660), gives the unique marked hyperbolic surface whose infinite grafting is the fixed pointed conformal surface $X$ and whose boundary lengths are $\ell_i$~\cite{C_Mondello2011}. Set $\lambda_i=\ell_i/(2\pi)$. For $\lambda_i>0$, denote its canonical cap by $D_i(\lambda)$ and its conformal radius in $z_i$ by $r_i(\lambda)>0$. For $\lambda_i=0$, the cap is the marked point and we set $r_i=0$.

We first work on a fixed face: the active index set $I=\{i:\lambda_i>0\}$ is fixed, and all other ends remain cusps. Put $x_i=\lambda_i^2$ for $i\in I$, and define
\[
 F_X:(0,\infty)^I\longrightarrow\mathbb R^I,
 \qquad F_X(x)_i=\log r_i(\sqrt{x}).
\]
The following argument concerns every stable $(g,n)$, including $(0,3)$, $(0,4)$ and $(1,1)$.

\begin{lemma}[Known regularity, with its application spelled out]
\label{C:A2b-regularity}
The caps and their normalized conformal coordinates vary real-analytically along each finite-dimensional real-analytic family in a fixed face of the length space. In particular $F_X$ is real-analytic. After identifying the infinitely grafted surface with $X\setminus\{p_i\}$, its Thurston metric varies differentiably off the moving seams. At a seam it is $C^{1,1}$, its first variation is continuous, and the one-sided derivatives used below exist.
\end{lemma}
\begin{proof}
Mondello's Theorem~5.4 states real-analyticity of the inverse on each finite topological stratum. We explain why it provides the local data needed here. A family of the resulting hyperbolic surfaces can be represented by varying Fenchel--Nielsen data; its boundary geodesics and Fermi collars vary real-analytically. With $s$ the signed hyperbolic distance from a boundary geodesic and $\theta$ its angle of period $2\pi$, the collar metric is
\[
 ds^2+\lambda_i^2\cosh^2(s)\,d\theta^2.
\]
The change $dy=ds/\cosh s$, $u=\lambda_i\theta$ is conformal and turns this into $\sec^2(y)(du^2+dy^2)$, with the seam at $y=0$. Continuing with $du^2+dy^2$ on $y\geq0$ gives the flat cylinder. For identifications with a fixed differentiable surface, keep $\theta$ and $y/\lambda_i$ fixed on each end and in a small collar. The metrics there change by a scalar conformal factor, so the Beltrami coefficient is zero on fixed neighborhoods of the punctures and seams. Interpolate these identifications over the remaining compact core. Its Beltrami coefficient depends real-analytically in $L^\infty$, while the seams stay fixed in this preliminary model. This is the local family construction underlying Mondello's proof in \S5.5, pp.~661--663.

Normalized solutions of the Beltrami equation depend differentiably, and real-analytically for analytic data, on these parameters: Ahlfors--Bers (1960), Theorem~10, Theorem~11 and its corollary, pp.~402--404~\cite{C_AhlforsBers1960}. Apply this to the marked conformal identifications with the fixed $X$. One may perform the construction in a finite conformal atlas, or on its normalized universal cover; uniqueness of the normalized solution makes the identifications compatible on overlaps. Because the conformal class is fixed, their targets are the same marked $X$. The preliminary end identifications just chosen are conformal, so the normalized Beltrami solutions are holomorphic on their fixed puncture and seam neighborhoods. Uniform parameter differentiability there implies differentiability of every holomorphic derivative by Cauchy's formula. On a cylinder the exponential
$w_i=\exp((-y+iu)/\lambda_i)$ is holomorphic and extends through the seam by the collar formula. Its normalization at $p_i$ fixes the harmless rotation. The parameter derivatives in the interior are holomorphic, and Cauchy's formula gives their derivatives on compact subsets. The same argument applies in the extended collar, so it includes both one-sided seam derivatives. Finally the two metric factors $1$ and $\sec^2(y)$ agree to first order at $y=0$ and have bounded second derivatives there. This gives the stated $C^{1,1}$ regularity and continuity of the first variation. Thus the cited parameter-dependence statements apply to the actual cap coordinates, not merely to abstract moduli points.
\end{proof}

For a tangent vector $v=(\dot x_i)_{i\in I}$, take a differentiable curve $x(t)$ with $x(0)=x$ and $\dot x(0)=v$. Write $g_t$ for its Thurston metric on the fixed punctured conformal surface, and let
\[
 \psi_v=\left.\frac{d}{dt}\log(g_t/g_0)\right|_{t=0}.
\]
This is a scalar function, independent of the holomorphic coordinate used to express the two conformal metrics. The base hyperbolic core is denoted $M$, and the base flat end belonging to $i$ by $Z_i$.

\begin{theorem}[Proved: the radius energy identity]
\label{C:A2b-energy}
For any two tangent vectors $v,w\in\mathbb R^I$, there are real harmonic functions $V_{i,v},V_{i,w}$ on the flat cylinders $Z_i$ such that
\begin{align}
4\pi\sum_{i\in I}v_i\,(dF_X)_x(w)_i
={}&\int_{X\setminus\{p_i\}}
 \bigl(\langle\nabla\psi_v,\nabla\psi_w\rangle_{g_0}
       +2\,1_M\psi_v\psi_w\bigr)\,dA_{g_0}\notag\\
 &+4\sum_{i\in I}\int_{Z_i}
       \langle\nabla V_{i,v},\nabla V_{i,w}\rangle\,dA.
\label{C:eq-radius-energy}
\end{align}
Every integral converges. The right-hand side is symmetric in $v,w$ and strictly positive when $v=w\ne0$. Consequently $DF_X(x)$ is symmetric positive definite at every $x$.
\end{theorem}
\begin{proof}
We give all normalizations and signs. Fix an active end and abbreviate its base length parameter by $\lambda$. Use flat coordinates
\[
 (u,y)\in(\mathbb R/(2\pi\lambda\mathbb Z))\times[0,\infty),
 \qquad w_0=\exp((-y+iu)/\lambda).
\]
The coordinate $w_t$ on the moving cap satisfies $|w_t|=1$ on its seam and has a simple zero at the fixed puncture. By Lemma~\ref{C:A2b-regularity}, it extends across the seam for small $t$. Its logarithmic variation
\[
 h_v=\left.\frac{d}{dt}\log w_t\right|_0
       =\frac{\dot w_t|_0}{w_0}
\]
is a single-valued holomorphic function on the base cap including its center and a collar of its boundary. Indeed $w_t$ vanishes at the fixed center, so the quotient has a removable singularity there. Define
\[
 a_v=\dot\lambda/\lambda,\qquad V_v=\lambda\operatorname{Re}h_v.
\]
Then $V_v$ is harmonic on $Z_i$. In the chosen local coordinate at the puncture,
$w_t=e^{i\alpha(t)}z_i/r_i(t)+O(z_i^2)$, and hence
\begin{equation}
 \lim_{y\to\infty}V_v(u,y)
   =-\lambda\,d(\log r_i)(v)=:b_{i,v}.
 \label{C:eq-radius-end}
\end{equation}
The convergence, with all derivatives, is exponential in $y/\lambda$. In particular the mean of $V_v$ on any horizontal circle is the constant $b_{i,v}$.

On the flat cap the metric is $g_t=\lambda(t)^2|dw_t/w_t|^2$. Differentiating this identity in any fixed holomorphic coordinate gives
\[
 \psi_v=2a_v+2\operatorname{Re}(w_0\partial_{w_0}h_v)
        =2a_v-2\partial_y V_v.
\]
Differentiating $\log|w_t|=0$ at its moving seam shows that the seam's normal velocity, positive from the hyperbolic core into the flat cylinder, is exactly $V_v(u,0)$. The base coordinate $u+iy$ extends to a hyperbolic collar $y<0$, where the metric is $\sec^2(y)|d(u+iy)|^2$. Thus the logarithmic metric factor and its first derivative agree across $y=0$, while the difference between its hyperbolic and flat second normal derivatives is $2$.

The curvature equation now gives, in the interiors,
\[
 \Delta_{g_0}\psi_v=2\psi_v\quad\hbox{on }M,
 \qquad \Delta_{g_0}\psi_v=0\quad\hbox{on each }Z_i,
\]
where $\Delta$ is the positive-coordinate Laplacian $\partial_u^2+\partial_y^2$ in a flat chart. Its interface condition is
\begin{equation}
 (\partial_y\psi_v)_{+}-(\partial_y\psi_v)_{-}=2V_v.
 \label{C:eq-radius-jump}
\end{equation}
One can obtain this by differentiating the matched first normal derivatives. Equivalently, the core is locally $\{y<tV_v(u,0)+o(t)\}$, so differentiating
$\Delta\phi_t=2e^{\phi_t}1_{M_t}$ gives the distributional term $2V_v\delta_{y=0}$. This also fixes the sign in~\eqref{C:eq-radius-jump}.

For clarity, all integrations can first be performed after truncating the flat cylinders at $y=T$. On a fixed cuspidal end, use its normalized local cusp coordinate $q_t$. The formula
\[
 g_t=\frac{|dq_t/q_t|^2}{(\log|q_t|)^2}
\]
and holomorphy of $\dot q_t/q_0$ show that, in $Y=-\log|q_0|$,
$\psi_v=O(Y^{-1})$ and $\partial_Y\psi_v=O(Y^{-2})$, up to exponentially decaying terms. These estimates make the cusp Dirichlet and potential integrals finite, and make its boundary terms tend to zero. On flat ends, $\psi_v\to2a_v$ and its derivatives decay exponentially; the corresponding boundary terms also tend to zero. These observations justify every following integration by parts and the passage to the untruncated surface.

Apply Green's identity to $\psi_v$ and the equation for $\psi_w$ on the core and all cylinders separately. Their outward normals at the seam are respectively $+\partial_y$ and $-\partial_y$. Using continuity of $\psi_v$ and~\eqref{C:eq-radius-jump}, we obtain
\begin{align}
 E_\psi(v,w)&:=\int
 \bigl(\langle\nabla\psi_v,\nabla\psi_w\rangle
          +2\,1_M\psi_v\psi_w\bigr)\,dA\notag\\
 &=-2\sum_{i\in I}\int_{\partial M_i}\psi_v V_{i,w}\,du.
 \label{C:eq-radius-green}
\end{align}
Harmonicity of $V_{i,v}$ similarly gives
\[
 E_{V_i}(v,w):=\int_{Z_i}\langle\nabla V_{i,v},\nabla V_{i,w}\rangle\,dA
 =-\int_{\partial M_i}(\partial_y V_{i,v})V_{i,w}\,du.
\]
Substitute $\psi_v=2a_{i,v}-2\partial_yV_{i,v}$ into~\eqref{C:eq-radius-green}. Since the circle has length $2\pi\lambda_i$ and the mean of $V_{i,w}$ is~\eqref{C:eq-radius-end}, this yields
\begin{align*}
 E_\psi(v,w)
 &=-4\sum_i a_{i,v}(2\pi\lambda_i)
                  (-\lambda_i\,d\log r_i(w))
       -4\sum_iE_{V_i}(v,w)\\
 &=4\pi\sum_i d(\lambda_i^2)(v)\,d\log r_i(w)
       -4\sum_iE_{V_i}(v,w).
\end{align*}
This is~\eqref{C:eq-radius-energy}.

The energy forms on its right are symmetric. If $v=w$ and the right-hand side vanishes, then $\psi_v=0$ on the nonempty hyperbolic core and $\nabla\psi_v=0$ everywhere. Each $V_{i,v}$ has zero Dirichlet energy and is constant. Continuity at the seam and $\psi_v=2a_{i,v}-2\partial_yV_{i,v}$ then give $a_{i,v}=0$ for every active $i$. Hence $v_i=2\lambda_i\dot\lambda_i=0$ for every $i$. Conversely zero variation gives zero energy. This proves strict positive definiteness.
\end{proof}

\begin{corollary}[Proved: global injectivity and a strictly convex potential]
\label{C:A2b-injectivity}
On every face of the nonnegative boundary-length space, the canonical cap system is determined by its conformal radii. On its positive coordinates the logarithmic radius map is a real-analytic diffeomorphism onto an open subset of $\mathbb R^I$, and is the gradient of a strictly convex real-analytic function of $x_i=(\ell_i/2\pi)^2$. In particular the full map of admissible canonical cap systems to their tangent-disc radii is injective. Changing the local coordinate norms translates the logarithmic radius map by a constant vector and does not change these assertions.
\end{corollary}
\begin{proof}
For distinct $x,x'$ in the convex positive orthant, put $v=x'-x$. The fundamental theorem of calculus and Theorem~\ref{C:A2b-energy} give
\[
 v\cdot(F_X(x')-F_X(x))
 =\int_0^1v^TDF_X(x+tv)v\,dt>0.
\]
Thus $F_X$ is injective. Its positive definite derivative is invertible, so the inverse function theorem makes it a local real-analytic diffeomorphism; injectivity then glues its local inverses into the claimed inverse on the open image. Symmetry of its derivative says that the one-form $\sum_i(F_X)_i\,dx_i$ is closed. Since the positive orthant is convex, it has a primitive $S_X$, whose Hessian is $DF_X$ and hence is positive definite. Explicitly, fixing $x_*$ one can take
\[
 S_X(x)=\int_0^1F_X(x_*+t(x-x_*))\cdot(x-x_*)\,dt.
\]
Its additive constant has no geometric significance.

If two radius vectors in the full nonnegative parameter space agree, their zero entries specify the same cuspidal indices; the remaining entries lie in the same face, where the preceding argument identifies their lengths. Mondello's uniqueness then identifies their cap systems on the fixed marked $X$. The all-cuspidal face consists of the single complete hyperbolic surface and requires no argument. Finally a coordinate change multiplies each radius by its fixed derivative modulus at $p_i$, so it adds a constant to $F_X$ and changes neither its derivative nor injectivity.
\end{proof}

\paragraph{The boundary argument needed for the image.}
The energy identity proves injectivity and local coordinates. To identify the open image, the next subsection supplies a separate support-gap estimate and properness argument. The image description is not inferred from positive definiteness alone.
